% TOP_PROBLEM: {"id":30007607,"problem_number":"LOCAL-30007607","title":"Nonbank systemic regulation","category_id":21,"category":{"id":21,"name":"economics","display_name":"Economics","description":"Open economic theory statements, published research agendas, and proposed scoped specifications; question provenance and historical priority are qualified per record.","slug":"economics","order_index":21,"created_at":"2026-10-08T00:00:00Z"},"status":"research_question","external_url":null,"legacy_ids":[],"published":false,"statement":"In the explicitly specified setting: Which leverage, redemption and margin rules best contain systemic risks in nonbank finance? The mathematical statement above fixes the target, assumptions and scope of this catalogue formulation.","economics":{"original_id":96,"field":"Banking, credit and financial stability","kind":"design","formalization_basis":"adapted_research_question","source_status":"explicit_open","priority_status":"earliest_found","priority_note":"This is the earliest explicit statement verified in this audit; first priority for the full catalogue formulation is not established.","earliest_verified_reference":{"authors":["Sirio Aramonte","Andreas Schrimpf","Hyun Song Shin"],"title":"Non-bank financial intermediaries and financial stability","year":2021,"url":"https://www.bis.org/publications/working-paper-972-non-bank-financial-intermediaries-and-financial-stability.pdf","doi":null,"locator":"Section 5, Open questions for policy and research, printed pp.38–39 (PDF pp.41–42), January 2022 revision","evidence_summary":"Explicitly asks optimal oversight and public-backstop design."},"current_reference":{"authors":["Bank of England"],"title":"Bank of England Agenda for Research, 2025–2028","year":2026,"url":"https://www.bankofengland.co.uk/research/bank-of-england-agenda-for-research","doi":null,"locator":"Prudential Architecture / Institutional investors, paragraph 3","evidence_summary":"Questions nonbank regulatory perimeter and new instruments."},"formalization_note":"Adapts an explicitly verified research-agenda passage. Mathematical objects, interventions, objectives and scope are proposed here; existing model-specific solutions are not relabeled unsolved."},"statusQualification":"Published question or research agenda verified at the cited locator. The mathematical specification is adapted for this catalogue and includes additional assumptions; source publication does not establish first-ever priority or certify this exact model remains unsolved. Adapts an explicitly verified research-agenda passage. Mathematical objects, interventions, objectives and scope are proposed here; existing model-specific solutions are not relabeled unsolved.","statusReviewedAt":"2026-10-08"}
% FIRST_POSTED: 2026-10-08
% ENTRY_KIND: statement_only
% !TeX program = lualatex
\documentclass[11pt]{article}
\usepackage{fontspec}
\usepackage[a4paper,margin=25mm]{geometry}
\usepackage{amsmath,amssymb,mathtools}
\usepackage{xurl}
\usepackage[hidelinks]{hyperref}
\newcommand{\sref}[2]{\hyperlink{source:#1:#2}{[S#2]}}
\setlength{\emergencystretch}{3em}
\begin{document}
% problemId:30007607
\section*{Nonbank systemic regulation}\label{30007607}
\subsection{Definitions and mathematical statement}
Consider nonbank intermediaries funding illiquid assets with redeemable claims, secured borrowing or derivatives. A feasible regulation \(a\) specifies leverage limits, redemption terms, swing pricing and initial or variation-margin rules. Investors choose leverage and liquidity buffers before shocks and trade or redeem after them; prices clear with explicitly modeled market impact. Let \(L(a,\omega)\) measure default, fire-sale and liquidity externalities in shock state \(\omega\), and \(C(a)\) routine intermediation and compliance costs. For a specified uncertainty set \(\mathcal P\), the design task is
\[a^*\in\arg\min_{a\in A}\left\{C(a)+\sup_{P\in\mathcal P}E_P[L(a,\omega)]\right\}.\]
The set \(A\) includes legal feasibility and operational implementability restrictions. Characterize how optimal tools depend on redemption mismatch, cross-institution exposures and endogenous margin spirals. Identify the information needed to estimate losses and costs, and test whether activity shifts to uncovered institutions when constraints tighten. Include changes in ordinary financing capacity, rather than treating regulation as costless. Report whether tools reduce individual-institution losses while increasing system-wide externalities. Existence of an optimal rule for a fixed shock law is a model result; robust implementation and empirical calibration remain separate parts of the research task.

\par\textbf{Formulation and scope.} Adapts an explicitly verified research-agenda passage. Mathematical objects, interventions, objectives and scope are proposed here; existing model-specific solutions are not relabeled unsolved.

\par\textbf{Open-question provenance.} An explicit question or research agenda was verified at the source locator. This does not by itself certify that every later version remains unresolved \sref{30007607}{1}.

\par\textbf{Historical priority.} This is the earliest explicit statement verified in this audit; first priority for the full catalogue formulation is not established.

\subsection{Short English statement}
In the explicitly specified setting: Which leverage, redemption and margin rules best contain systemic risks in nonbank finance? The mathematical statement above fixes the target, assumptions and scope of this catalogue formulation.

\subsection{Sources}
\begin{itemize}\raggedright
\item[\textnormal{[S1]}]\hypertarget{source:30007607:1}{} Sirio Aramonte, Andreas Schrimpf, Hyun Song Shin. 2021. Non-bank financial intermediaries and financial stability. Section 5, Open questions for policy and research, printed pp.38–39 (PDF pp.41–42), January 2022 revision.
\url{https://www.bis.org/publications/working-paper-972-non-bank-financial-intermediaries-and-financial-stability.pdf}.
{\small\textit{Earliest verified question/agenda reference; historical priority qualified below. Explicitly asks optimal oversight and public-backstop design.}}

\item[\textnormal{[S2]}]\hypertarget{source:30007607:2}{} Bank of England. 2026. Bank of England Agenda for Research, 2025–2028. Prudential Architecture / Institutional investors, paragraph 3.
\url{https://www.bankofengland.co.uk/research/bank-of-england-agenda-for-research}.
{\small\textit{Supporting or current literature; see evidence qualification. Questions nonbank regulatory perimeter and new instruments.}}
\end{itemize}
\end{document}
